% Actual class: 10
% Slide deck: Part 4 (Process Synchronization)
% Exact scope: pp. 37--54 (substantive content pp. 38--54)
% Video supplement: Part 4.7 and Part 4.8 complete
% Human verified: Yes

\section{第 10 节课：Semaphore 设计与经典同步问题}
\label{lecture:SC2005-L10}

\lecturemeta
  {SC2005-L10}
  {2026-09-11}
  {Part 4 (Process Synchronization)}
  {主讲义 pp.~37--54（实质 pp.~38--54）}
  {Part 4.7 Semaphore Best Practices 与 Part 4.8 Classical Synchronization Problems 完整}
  {已确认（2026-09-11）}

\subsection{本讲主线}

Semaphore（信号量）题的关键不是背代码，而是先区分两类约束：shared data（共享数据）需要 mutual exclusion（互斥），resource availability / execution order（资源可用性／执行顺序）需要 counting 或 binary semaphore 表示条件。错误的获取顺序可能造成 deadlock（死锁），不公平的唤醒策略还可能造成 starvation（饥饿）。

\lecturetopic{SC2005-L10-T01}{Semaphore 的系统化设计}

\begin{enumerate}
  \item 找出所有 shared variables 与 shared resources。
  \item 对需要原子访问的 shared data，用初值为 $1$ 的 binary semaphore（通常命名为 \texttt{mutex}）包围 critical section。
  \item 对每类可计数资源各设一个 semaphore，初值等于资源实例数；取得／消耗资源前执行 \texttt{Wait}，释放／产生资源后执行 \texttt{Signal}。
\end{enumerate}

Semaphore 也可表达 happens-before（先发生）关系。若要求 $A$ 必须先于 $B$：
\[
\begin{array}{ll}
P_i: & A;\ \texttt{Signal(flag);}\qquad (\texttt{flag}=0),\\
P_k: & \texttt{Wait(flag);}\ B.
\end{array}
\]
$B$ 在取得由 $A$ 产生的唯一 permit 前无法执行。

\textbf{对应独立检验例：}\relatedexercise{SC2005-L10-E01}{用 Semaphore 保证执行顺序}。

\lecturetopic{SC2005-L10-T02}{Deadlock 与 Starvation}

\begin{itemize}
  \item \textbf{Deadlock（死锁）：}一组 processes 形成 circular wait（循环等待），每个都在等待只能由组内另一个 process 触发的事件，因此整个集合都不能继续。
  \item \textbf{Starvation（饥饿）：}系统仍在运行，但某个 process 因调度、优先级或不公平的 semaphore queue 长期得不到资源。
\end{itemize}

例如 $S=Q=1$：$P_0$ 先取得 $S$ 后等待 $Q$，$P_1$ 先取得 $Q$ 后等待 $S$，二者便互相等待。统一 lock ordering（锁顺序）可以破坏 circular wait。FIFO waiting queue 有助于 bounded waiting，但仍需要资源最终被释放和 scheduler 公平等假设。

\textbf{对应独立检验例：}\relatedexercise{SC2005-L10-E02}{两个 Semaphore 的 Circular Wait}。

\lecturetopic{SC2005-L10-T03}{Bounded-Buffer Producer--Consumer}

容量为 $n$ 的 circular buffer（循环缓冲区）包含 shared data \texttt{buffer}, \texttt{in}, \texttt{out}，并使用
\[
  \texttt{mutex}=1,\qquad \texttt{empty}=n,\qquad \texttt{full}=0.
\]
其中 \texttt{mutex} 保护 buffer 与索引；\texttt{empty}/\texttt{full} 分别记录 empty slots（空槽）与 full slots（已占槽），因此不再需要另设共享 \texttt{counter}。

\begin{center}
\small
\begin{tabularx}{0.94\textwidth}{@{}XX@{}}
  \toprule
  Producer（生产者） & Consumer（消费者） \\
  \midrule
  produce item & \texttt{Wait(full)} \\
  \texttt{Wait(empty)} & \texttt{Wait(mutex)} \\
  \texttt{Wait(mutex)} & remove item; update \texttt{out} \\
  insert item; update \texttt{in} & \texttt{Signal(mutex)} \\
  \texttt{Signal(mutex)} & \texttt{Signal(empty)} \\
  \texttt{Signal(full)} & consume item \\
  \bottomrule
\end{tabularx}
\end{center}

\begin{figure}[htbp]
  \centering
  \resizebox{0.96\textwidth}{!}{\begin{tikzpicture}[
  >=Latex,
  actor/.style={draw=noteBlue,very thick,rounded corners=3pt,fill=blue!6,minimum width=2.15cm,minimum height=9mm,align=center,font=\small},
  resource/.style={draw=ntuRed,very thick,rounded corners=3pt,fill=red!5,minimum width=2.15cm,minimum height=9mm,align=center,font=\small},
  mutex/.style={draw=verifyOrange,very thick,rounded corners=3pt,fill=orange!7,minimum width=2.15cm,minimum height=9mm,align=center,font=\small},
  access/.style={draw=green!45!black,very thick,rounded corners=3pt,fill=green!7,minimum width=2.15cm,minimum height=9mm,align=center,font=\small},
  flow/.style={-{Latex[length=2mm]},thick,draw=black!68}
]
  \node[actor] (p0) at (0,1.05) {Producer};
  \node[resource] (p1) at (2.65,1.05) {Wait(empty)};
  \node[mutex] (p2) at (5.3,1.05) {Wait(mutex)};
  \node[access] (p3) at (7.95,1.05) {insert item};
  \node[mutex] (p4) at (10.6,1.05) {Signal(mutex)};
  \node[resource] (p5) at (13.25,1.05) {Signal(full)};

  \node[actor] (c0) at (0,-1.05) {Consumer};
  \node[resource] (c1) at (2.65,-1.05) {Wait(full)};
  \node[mutex] (c2) at (5.3,-1.05) {Wait(mutex)};
  \node[access] (c3) at (7.95,-1.05) {remove item};
  \node[mutex] (c4) at (10.6,-1.05) {Signal(mutex)};
  \node[resource] (c5) at (13.25,-1.05) {Signal(empty)};

  \foreach \a/\b in {p0/p1,p1/p2,p2/p3,p3/p4,p4/p5,c0/c1,c1/c2,c2/c3,c3/c4,c4/c5}
    \draw[flow] (\a) -- (\b);

  \node[font=\scriptsize,text=ntuRed,anchor=east] at (13.25,0) {resource count};
  \node[font=\scriptsize,text=verifyOrange] at (5.3,0) {protect shared data};
  \node[font=\scriptsize,text=green!45!black] at (7.95,0) {buffer access};
\end{tikzpicture}
}
  \caption{Bounded buffer 中三个 semaphores 的职责。Resource semaphore 必须先于 \texttt{mutex} 获取，避免持锁等待。}
  \label{fig:sc2005-l10-bounded-buffer}
\end{figure}

\keyidea{先等待所需资源，再进入保护 shared data 的 critical section；不要持有 \texttt{mutex} 去等待只能由另一个需要该 \texttt{mutex} 的 process 创造的条件。}

若 buffer 已满，producer 错误地先 \texttt{Wait(mutex)}、再 \texttt{Wait(empty)}，它会持有 \texttt{mutex} 阻塞；consumer 虽通过 \texttt{Wait(full)}，却无法取得 \texttt{mutex} 来移除 item 并 \texttt{Signal(empty)}，于是 deadlock。

\textbf{对应独立检验例：}\relatedexercise{SC2005-L10-E03}{Bounded Buffer 的错误 Wait 顺序}。

\lecturetopic{SC2005-L10-T04}{Dining Philosophers}

五位 philosophers 与五根 chopsticks 构成环。若每人都先取得左边 chopstick、再等待右边 chopstick，可能出现五人各持一根并等待下一根的 circular wait。

\begin{figure}[htbp]
  \centering
  \begin{tikzpicture}[
  >=Latex,
  philosopher/.style={circle,draw=noteBlue,very thick,fill=blue!6,minimum size=10mm,font=\small},
  chopstick/.style={circle,draw=ntuRed,very thick,fill=red!6,minimum size=7mm,font=\scriptsize},
  wait/.style={-{Latex[length=2mm]},thick,draw=black!68}
]
  \foreach \i in {0,...,4} {
    \pgfmathsetmacro{\a}{90-72*\i}
    \pgfmathsetmacro{\b}{54-72*\i}
    \node[philosopher] (p\i) at (\a:2.25cm) {$P_{\i}$};
    \node[chopstick] (c\i) at (\b:1.35cm) {$C_{\i}$};
  }
  \foreach \i in {0,...,4} {
    \pgfmathtruncatemacro{\j}{mod(\i+1,5)}
    \draw[wait] (p\i) to[bend left=10] (c\i);
    \draw[wait,dashed] (c\i) to[bend left=10] (p\j);
  }
  \node[align=center,font=\scriptsize] at (0,-2.95cm)
    {实线：等待下一根\qquad 虚线：已被相邻者持有};
\end{tikzpicture}

  \caption{所有 philosophers 同时持有左侧 chopstick 时形成 circular wait。}
  \label{fig:sc2005-l10-dining-philosophers}
\end{figure}

常见 deadlock prevention（死锁预防）包括：
\begin{itemize}
  \item 最多允许四位 philosophers 同时尝试取 chopsticks；至少留下一根空闲 chopstick，使相邻者能够取得第二根、完成并释放资源。
  \item 将“同时取得两根 chopsticks”做成 atomic operation。
  \item 采用 asymmetric ordering（非对称顺序）或全局资源编号顺序，例如奇偶 philosopher 以相反顺序取用。
\end{itemize}
这些方案破坏 deadlock 的必要条件，但不自动保证没有 starvation；公平性仍取决于 queue 与 scheduling policy。

\textbf{对应独立检验例：}\relatedexercise{SC2005-L10-E04}{限制四位 Philosophers 的作用}。

\lecturetopic{SC2005-L10-T05}{Readers--Writers Problem}

Readers 可同时读取 shared database；writer 写入时必须独占。设
\[
  \texttt{readcount}=0,\qquad \texttt{mutex}=1,\qquad \texttt{wrt}=1.
\]
\texttt{mutex} 只保护 \texttt{readcount}；\texttt{wrt} 控制 database 的读者群体或单个 writer。

\begin{center}
\small
\begin{tabularx}{0.94\textwidth}{@{}XX@{}}
  \toprule
  Reader（读者） & Writer（写者） \\
  \midrule
  \texttt{Wait(mutex); readcount++;} & \texttt{Wait(wrt)} \\
  \texttt{if (readcount == 1) Wait(wrt);} & write database \\
  \texttt{Signal(mutex);} & \texttt{Signal(wrt)} \\
  read database & \\
  \texttt{Wait(mutex); readcount--;} & \\
  \texttt{if (readcount == 0) Signal(wrt);} & \\
  \texttt{Signal(mutex);} & \\
  \bottomrule
\end{tabularx}
\end{center}

First reader（首位读者）替整个 reader group 取得 \texttt{wrt}，last reader（末位读者）释放它；中间读者只更新 \texttt{readcount}。该 reader-preference（读者优先）版本允许新 readers 持续加入，因此 writer 可能 starvation。

\textbf{对应独立检验例：}\relatedexercise{SC2005-L10-E05}{Readers--Writers 状态轨迹}。

\lecturetopic{SC2005-L10-T06}{Monitor：不考的延伸}

Monitor（管程）是封装 shared data、operations 与 synchronization 的高层语言结构，同一时刻只允许一个 process 在 monitor 内 active。它减少手工配对 \texttt{Wait}/\texttt{Signal} 的错误；本课程将此部分标为 not examinable（不考），只需认识概念。

\textbf{一分钟回顾。}\texttt{empty/full} 管理资源，\texttt{mutex} 保护数据；统一 resource ordering 可避免 circular wait。Deadlock freedom 与 starvation freedom 必须分开检查。
